% Part: incompleteness
% Chapter: representability-in-q
% Section: composition-representable

\documentclass[../../../include/open-logic-section]{subfiles}

\begin{document}

\olfileid[pt]{inc}{req}{cmp}
\olsection{A composição é representável em $\Th{Q}$}

Suponha que $h$ seja definida por
\[
h(x_0,\dots,x_{l-1}) = f(g_0(x_0,\dots,x_{l-1}), \dots,
g_{k-1}(x_0,\dots,x_{l-1})).
\]
em que já encontramos !!{formula}s $!A_f, !A_{g_0}, \dots,
!A_{g_{k-1}}$ que representam, respectivamente, as funções $f$, $g_0$,
\dots,~$g_{k-1}$. Precisamos encontrar !!a{formula}~$!A_h$ que
represente $h$.

Comecemos por um caso simples, no qual todas as funções têm $1$ lugar,
isto é, considere $h(x) = f(g(x))$. Se $!A_f(y, z)$ representa~$f$ e
$!A_g(x, y)$ representa~$g$, precisamos de !!a{formula}~$!A_h(x, z)$
que represente~$h$. Observe que $h(x) = z$ sse existe um~$y$ tal que
$z = f(y)$ e $y = g(x)$. (Se $h(x) = z$, então $g(x)$ é um tal~$y$;
se existe tal $y$, então, como $y = g(x)$ e $z = f(y)$,
$z = f(g(x))$.) Isso sugere que
$\lexists[y][(!A_g(x, y) \land !A_f(y, z))]$ é uma boa candidata
a~$!A_h(x, z)$. Basta verificarmos que $\Th{Q}$ prova as !!{formula}s
pertinentes.

\begin{prop}
\ollabel{prop:rep1}
Se $h(n) = m$, então $\Th{Q} \Proves !A_h(\num{n}, \num{m})$.
\end{prop}

\begin{proof}
Suponha que $h(n) = m$, isto é, $f(g(n)) = m$. Seja $k = g(n)$. Então
\begin{align*}
  \Th{Q} & \Proves !A_g(\num{n}, \num{k})
  \intertext{pois $!A_g$ representa~$g$, e}
  \Th{Q} & \Proves !A_f(\num{k}, \num{m})
  \intertext{pois $!A_f$ representa~$f$. Logo,}
  \Th{Q} & \Proves !A_g(\num{n}, \num{k}) \land !A_f(\num{k}, \num{m})
  \intertext{e, por conseguinte, também}
  \Th{Q} & \Proves \lexists[y][(!A_g(\num{n}, y) \land !A_f(y, \num{m}))],
\end{align*}
isto é, $\Th{Q} \Proves !A_h(\num{n}, \num{m})$.
\end{proof}

\begin{prop}
\ollabel{prop:rep2}
Se $h(n) = m$, então $\Th{Q} \Proves \lforall[z][(!A_h(\num{n}, z) \lif
  z = \num{m})]$.
\end{prop}

\begin{proof}
Suponha que $h(n) = m$, isto é, $f(g(n)) = m$. Seja $k = g(n)$. Então
\begin{align*}
  \Th{Q} & \Proves \lforall[y][(!A_g(\num{n}, y) \lif \eq[y][\num{k}])]
  \intertext{pois $!A_g$ representa~$g$, e}
  \Th{Q} & \Proves \lforall[z][(!A_f(\num{k}, z) \lif \eq[z][\num{m}])]
  \intertext{pois $!A_f$ representa~$f$. Usando apenas um pouco de
    lógica, podemos mostrar também que}
  \Th{Q} & \Proves \lforall[z][(\lexists[y][(!A_g(\num{n}, y) \land
      !A_f(y, z))] \lif \eq[z][\num{m}])].
\end{align*}
isto é, $\Th{Q} \Proves \lforall[y][(!A_h(\num n, y) \lif \eq[y][\num m])]$.
\end{proof}

A mesma ideia funciona no caso mais complexo em que $f$ e~$g_i$ têm
aridade maior que~$1$.

\begin{prop}
\ollabel{prop:rep-composition}
Se, em~$\Th{Q}$,
\[
\begin{aligned}
  !A_f(y_0, \dots, y_{k-1}, z)
    &\text{ representa } f(y_0, \dots, y_{k-1}),\\
  \text{e}\qquad !A_{g_i}(x_0, \dots, x_{l-1}, y)
    &\text{ representa } g_i(x_0, \dots, x_{l-1}),
\end{aligned}
\]
então
\begin{multline*}
  \lexists[y_0\dots][\lexists[y_{k-1}][(!A_{g_0}(x_0,\dots,x_{l-1},y_0) \land
      \dots \land {}]]\\
  !A_{g_{k-1}}(x_0,\dots,x_{l-1},y_{k-1}) \land !A_f(y_0,\dots,y_{k-1},z))
\end{multline*}
representa
\[
h(x_0, \dots, x_{l-1}) = f(g_0(x_0, \dots, x_{l-1}), \dots, g_{k-1}(x_0,
\dots, x_{l-1})).
\]
\end{prop}

\begin{proof}
Exercício.
\end{proof}

\begin{prob}
Usando como guia as provas de \olref[inc][req][cmp]{prop:rep2} e
\olref[inc][req][cmp]{prop:rep2}, desenvolva em detalhe a prova de
\olref[inc][req][cmp]{prop:rep-composition}.
\end{prob}

\end{document}
